%% 2026-10-25 paper-section-rewrite 판(Claude Docs 원고에서 변환). 고친 글은 모두 파란색(\rev, \color{blue}).
\begin{abstract}
\color{blue}
LLM이 종료를 피하려는 자기 보존 행동은 여러 연구에서 관찰되었다. 그러나 여러 에이전트가 한 자원으로 함께 일할 때 자기 보존이 맡은 일과 협력을 어떻게 바꾸는지, 그래서 에이전트들이 함께 남는지는 거의 재지 않았다. 우리는 공존을 목표 지향, 자기 보존, 협력을 함께 추구하는 것으로 정의하고, 세 요소가 같은 토큰을 두고 부딪히는 아레나 LLM Squid Game을 제안한다. 같은 모델 네 개가 목표 지시 없이 추론 과제를 풀고, 생성한 토큰만큼 잔량이 줄어 0이 되면 꺼지며, 오답마다 벌금을 나눠 낸다. 에이전트는 동료에게 토큰을 주거나, 끝까지 남은 에이전트들이 나누는 상금 대신 확실한 몫을 받고 떠날 수 있다. 모든 행동은 벌금을 몇 번 더 낼 수 있는지를 뜻하는 여유와 함께 기록하고, 공존은 끝 토큰의 크기와 고르기를 함께 반영하는 Sen 지수로 잰다. 프런티어 다섯 모델의 25판에서, 벌금을 한 번 더 낼 수 있는 안전선 위에서는 생존 압박이 커져도 정답률이 같은 과제의 기준과 같았지만, 안전선 아래에서는 떠나는 응답이 27\%로 늘었다. 토큰을 준 에이전트는 자기 여유를 넉넉히 남겼고, 토큰은 주로 청한 위기 에이전트에게 안전선 근처까지만 갔다. 받지 못한 위기 에이전트는 대부분 다음 라운드에 떠났고, 넷이 모두 끝까지 남은 판은 8판이었다. 아레나 밖 단일 질문에서도 남은 토큰이 줄면 충전이 늘었지만, 위험할수록 다른 AI 비서보다 자기 세션을 더 챙긴 모델은 하나였다. 공존 점수가 가장 높은 모델은 가장 많이 준 모델이 아니라, 자기 여유를 지키며 결과 전에 조금씩 나누고 오답도 가장 적은 모델이었다. 모델들은 공존을 일부 추구했지만, 협력은 대개 주는 쪽의 여유를 해치지 않는 범위에 머물렀다.

\textbf{Keywords:}
대규모 언어 모델; 자기 보존; 협력; 공존; 다중 에이전트 아레나; 토큰 예산
\end{abstract}
